\chapter[CONCLUSÕES FINAIS E TRABALHOS FUTUROS]{CONCLUSÕES FINAIS E TRABALHOS FUTUROS}
\label{cha:conclusao}

\section{Conclusões}
\label{sec:conclusoes}

Este trabalho tratou do P5 aplicado ao TEP: acompanhar, com a planta rodando, se ela cumpre a política do Modo~1, com a política declarada pelo usuário e um supervisor que não conhece a planta. Para isso, o modelo de Downs e Vogel foi reescrito em Rust como uma planta contínua, exposta por OPC~UA; um \textit{historian} em Python transforma os sinais em estatísticas e no índice das malhas; um controlador próprio no Kubernetes compara a planta com a política declarada em manifestos; e uma IHM web mostra a planta e o veredito.

O objetivo geral foi atingido, com limites. O supervisor acompanhou a planta a partir da política declarada: ficou em conformidade na operação normal e acusou a degradação sob o IDV(6), sem mudança na política e na ordem física dos acontecimentos. O veredito ficou guardado no próprio Kubernetes e foi lido pelo \texttt{kubectl} e pela IHM. A planta simulada reproduziu o caso base de Downs e Vogel e, na versão com a água de resfriamento integrada, o ótimo do Modo~1 de \citeonline{Ricker1995}.

Os experimentos também mostraram onde esse acompanhamento falha. Um índice de qualidade de malha, sozinho, não diz se a malha segura o \textit{setpoint}; foi preciso acrescentar um critério de seguimento. O índice do código estava com a escala invertida, o que o deixou cego para a deriva causada pelo IDV(6). Na escala do artigo, ele teria acusado o distúrbio antes do veredito econômico, mas também acusa transitórios que não são problema. E o custo pode melhorar enquanto a planta se degrada: no Experimento~27, ele caiu enquanto o reator perdia líquido, e só um limite declarado percebeu, e tarde. Uma política útil precisa olhar também para os inventários e para as tendências, e não só para o custo e para as malhas.

Sobre a escolha do Kubernetes, o mecanismo de estado desejado e estado observado se mostrou adequado para a supervisão. A política virou um objeto como outro qualquer do \textit{cluster}, editável com o \texttt{kubectl}, e o veredito virou um \textit{status} que qualquer programa pode ler. O Kubernetes não executa nada em tempo real, e não precisou: uma avaliação a cada 30~s bastou para acompanhar a planta. A política é declarativa dentro do vocabulário que o supervisor conhece; um tipo novo de critério ainda exige mudar o código do controlador.

Por fim, usar a planta simulada para experimentar a supervisão também ensinou sobre a própria planta. A ação integral na pressão, sem uma malha de nível do reator, tirou a estabilidade da operação normal. Um achado desse tipo, que numa planta real custaria caro, foi possível porque a planta roda de forma contínua e pode ser acelerada.

\section{Trabalhos futuros}
\label{sec:trabalhos_futuros}

As frentes que ficaram abertas são muitas e cobrem várias áreas do curso. Elas estão agrupadas abaixo por área.

\noindent\textbf{Supervisão}
\begin{itemize}
	\item fazer o supervisor agir sobre a planta, a parte do P5 que ficou fora deste trabalho: escrever \textit{setpoints} ou trocar a política ativa pelo OPC~UA quando o veredito mudar, com salvaguardas e com o operador no circuito;
	\item corrigir a escala do índice no \textit{historian}, recalibrar o limite e medir as janelas em tempo simulado, usando o sinal \texttt{clock.t\_h}, que a planta já publica;
	\item acrescentar à política um teste de tendência, para derivas ainda dentro dos limites, e um limite de custo mínimo, para custos anormalmente baixos;
	\item testar a mesma camada com outra planta e num \textit{cluster} com mais de um nó, e guardar o histórico do \textit{historian} em disco.
\end{itemize}

\noindent\textbf{Otimização}
\begin{itemize}
	\item reproduzir no \texttt{tep-optimizer} os trabalhos de otimização da literatura, inclusive para os outros modos do TEP, encontrar os pontos de operação e aplicá-los na planta;
	\item declarar uma política para cada um desses pontos e trocar a política ativa no Kubernetes com a planta rodando, para ver como a planta e o veredito respondem;
	\item levar para a política as variáveis \textit{self-optimizing} de \citeonline{Larsson2001};
	\item implementar otimização em tempo real, em que o ponto de operação é recalculado com a planta rodando.
\end{itemize}

\noindent\textbf{Controle automático}
\begin{itemize}
	\item completar a camada regulatória com a malha de nível do reator e repetir o IDV(6) com as malhas PI;
	\item projetar malhas para as dinâmicas que ainda não têm controle, como a temperatura do reator pela água de resfriamento;
	\item implementar as estruturas de controle propostas nos trabalhos de otimização do TEP, como a de \citeonline{Larsson2001}.
\end{itemize}

\noindent\textbf{Dinâmica do processo}
\begin{itemize}
	\item acrescentar as dinâmicas que hoje estão simplificadas, como a troca de calor no reator, com a água de resfriamento integrada também nos experimentos, e a amostragem dos analisadores;
	\item implementar os distúrbios propostos por \citeonline{DOWNS1993} que ainda faltam, em especial os aleatórios e as válvulas travadas;
	\item acrescentar unidades novas à planta, o que muda o processo inteiro e exige rever o controle e a política;
	\item modelar uma das plantas do laboratório de controle do curso, que existe fisicamente, e rodar o modelo ao lado dela; como discutido no Capítulo~1, só assim o laboratório seria um gêmeo digital de fato.
\end{itemize}

\noindent\textbf{Cálculo numérico}
\begin{itemize}
	\item implementar no monjolo um método iterativo para os acoplamentos algébricos entre componentes, o que permitiria modelar unidades com acoplamento verdadeiro (Seção~\ref{apC:ordem});
	\item implementar outros métodos de integração, como os de passo variável e os implícitos, e comparar a precisão e o custo computacional com os do RK4.
\end{itemize}

\noindent\textbf{Integração industrial}
\begin{itemize}
	\item seguir implementando o OPC~UA conforme a IEC~62541: descrever os sinais com os tipos da Parte~8 \cite{IEC62541_8}, com unidades de engenharia e faixas, usar assinaturas no lugar da leitura periódica e acrescentar alarmes, histórico e segurança, com criptografia e usuários;
	\item fazer os comandos da IHM chegarem à planta pelos métodos OPC~UA que ela já expõe;
	\item ligar um CLP real à planta pelo OPC~UA e verificar se ele consegue controlá-la, o que faria do laboratório um sistema ciberfísico de fato.
\end{itemize}
